% ============================================================
%  Etymology, security definitions, threat models, pillars
% ============================================================

\section{From ``hidden writing'' to modern cryptography}

Let's start from the name. The word \textbf{cryptography} comes from ancient Greek:
\begin{itemize}
  \item \textit{kryptós} $\to$ ``hidden, secret'' (from the verb \textit{krýptein}, ``to hide'');
  \item \textit{gráphein} $\to$ ``to write''.
\end{itemize}
Literally, cryptography means \textbf{``hidden writing''}. The same root gives \textbf{cryptanalysis} (from \textit{analýein}, ``to untie'': the art of breaking ciphers) and \textbf{cryptology} (from \textit{lógos}, ``study'': the discipline that includes both).

\subsection{Cryptography versus steganography}

The name already tells us \textit{what} is hidden, and it helps to compare cryptography with a similar word:
\begin{itemize}
  \item \textbf{Cryptography} $\to$ hides the \textbf{meaning} of the message: everybody can see that a message exists, but nobody can read it.
  \item \textbf{Steganography} (from \textit{steganós}, ``covered'') $\to$ hides the \textbf{existence} of the message: for example, a text hidden inside the pixels of an image, where nobody suspects anything.
\end{itemize}
This difference is more important than it looks: \textbf{in cryptography we never rely on the adversary not noticing the message}. We always assume that she sees the ciphertext and knows the algorithm, and that the only secret is the key. This is Kerckhoffs's principle, which we will meet in \Cref{sec:ske}.

\subsection{From an art to a science}

For centuries, from Caesar's cipher to the Enigma machine, cryptography was exactly what its name says: hidden writing, with a single goal, \textbf{confidentiality}. And it was an \textit{art}: someone invented a cipher, someone else broke it, and the cycle started again. There was no way to say in advance whether a cipher was good: ``secure'' simply meant ``nobody has broken it yet''.

Modern cryptography changed in \textbf{two ways}:
\begin{itemize}
  \item \textbf{The goals became broader} $\to$ not only hiding the message (\textit{confidentiality}), but also detecting modifications (\textit{integrity}) and verifying who sent it (\textit{authenticity}). A digital signature hides nothing at all, and yet it is cryptography: the name ``hidden writing'' now tells only one third of the story. These three goals are the \textbf{pillars of security} (\Cref{sec:pillars}).
  \item \textbf{It became a science} $\to$ instead of hoping that nobody breaks a cipher, we first \textbf{define} mathematically what ``secure'' means, and then \textbf{prove} that a scheme satisfies the definition.
\end{itemize}

\begin{frameddefn}{Modern cryptography}
  Modern cryptography is the scientific study of techniques that protect information and communication \textbf{in the presence of adversaries}, based on \textbf{precise definitions of security} and on \textbf{mathematical proofs}.
\end{frameddefn}

Both changes lead to the same question, which is the starting point of the whole course: \textbf{what does ``secure'' mean, exactly?}

\section{What does ``secure'' mean?}

The question looks trivial: \textit{when is a system secure?} In everyday language ``secure'' is a feeling: something is secure if nobody has managed to break it yet. Modern cryptography cannot work like this. For us, ``secure'' is a \textbf{precise mathematical statement}: first we \textbf{define} what security means, then we \textbf{prove} that our scheme satisfies the definition. A scheme without a security definition cannot be called secure, simply because there is nothing to prove.

To write a security definition we must answer \textbf{two questions}:
\begin{itemize}
  \item \textbf{What must we protect?} (the \textbf{security goal}) $\to$ what the attacker \textbf{must not manage to do}. For example:
  \begin{itemize}
    \item learn information about the message $\to$ this is \textit{confidentiality};
    \item modify or forge a message without being detected $\to$ these are \textit{integrity} and \textit{authenticity}.
  \end{itemize}
  \item \textbf{Who are we protecting it from?} (the \textbf{threat model}) $\to$ what \textbf{resources} the attacker has, and what we allow the attacker to \textbf{see and do}.
\end{itemize}

Note the wording: the attacker ``\underline{must not manage to}'' do something. It is not an observation (``the attacker cannot''), it is a \textbf{requirement}, and it is exactly what the security proof will have to show.

Once we have fixed these two things, every security definition is built following the same template.

\begin{framedobs}[label={obs:template}]{The template of every security definition (informal)}
  A scheme $\Pi$ is \textbf{secure}, with respect to a security goal and a threat model, if no adversary allowed by the threat model can do better than a random guess, \textbf{except with a negligible (very small) probability}:
  \[
    \forall \Adv \in \text{threat model}: \quad \big|\Pr[\Adv \text{ wins } \mathsf{G}_{\Pi}] - p^*\big| \leq \varepsilon
  \]
  \begin{itemize}
    \item $\mathsf{G}_{\Pi}$ $\to$ the \textbf{security game}: a formal ``challenge'' between the scheme and the adversary;
    \item $p^*$ $\to$ the probability of winning by \textbf{trivially guessing} (for example $\frac{1}{2}$ when the adversary has to guess a bit);
    \item $\varepsilon$ $\to$ the \textbf{advantage} we tolerate: $0$ for information-theoretic security, a negligible function $\negl(n)$ of the security parameter $n$ for computational security.
  \end{itemize}
  This is \textbf{not} a rigorous definition: ``wins the game'' and ``threat model'' become precise only when we fix a \textbf{specific game} and a \textbf{specific class of adversaries}. The rigorous versions are in \Cref{sec:rigorous}.
\end{framedobs}

Be careful not to shorten the template to ``no adversary wins, except with small probability'': that would be wrong. In the guessing game below the adversary wins half of the times just by tossing a coin, and $\frac{1}{2}$ is not a small probability. What must be small is the probability of winning \textbf{more than} by guessing.

\subsection{The security game, explained with an example}

What is this ``game''? It is simply a \textbf{quiz} that tests the adversary. Imagine that a general sends a single bit: \textbf{0} = ``attack'', \textbf{1} = ``retreat''.
\begin{enumerate}
  \item Eve knows that the message is either ``attack'' or ``retreat''.
  \item Secretly, I toss a coin, pick one of the two messages, encrypt it and give Eve \textbf{only the ciphertext}.
  \item Eve looks at the ciphertext and must say which message it was.
\end{enumerate}
We repeat the quiz many times and count how often Eve guesses correctly:
\begin{itemize}
  \item \textbf{No encryption} $\to$ Eve reads the message and guesses 100\% of the times $\to$ the scheme is broken;
  \item \textbf{A flawed scheme} $\to$ Eve guesses, for example, 75\% of the times $\to$ she learned something, the scheme is broken;
  \item \textbf{One-Time Pad} $\to$ Eve guesses exactly 50\% of the times $\to$ the scheme is secure.
\end{itemize}

Eve \textbf{always} gets 50\%, even without looking at the ciphertext: it is enough to toss a coin. So the real question is: \textbf{does looking at the ciphertext help Eve do better than a coin?} If not, the scheme is secure.

In the official version of the game it is the \textit{adversary} who chooses the two messages. This way she picks the pair that is most favourable to her, and if the scheme resists even in that case, it resists for every pair of messages.

\subsection{The threat model}

There is no such thing as security ``in absolute terms''. A safe can be perfectly secure against a thief with a screwdriver and useless against a thief with a blowtorch. So, before proving anything, we must describe the adversary precisely, along \textbf{two axes}.

\paragraph{First axis: computational power} (how much can the adversary compute?)
\begin{itemize}
  \item \textbf{Unbounded} $\to$ infinite time and memory: the adversary can try every key, every message, any computation. It is the \textbf{strongest possible} adversary: if a scheme resists it, it resists anyone. This is the model of perfect secrecy and of the One-Time Pad.
  \item \textbf{Polynomial-time} $\to$ the adversary is an \textbf{efficient algorithm}: a real computer, even a very powerful one, but it cannot perform, for example, $2^{128}$ operations. It is the \textbf{realistic} model, used in almost all the rest of the course.
\end{itemize}

\paragraph{Second axis: capabilities} (what can the adversary see and do?)

\begin{table}[H]
  \centering
  \small
  \begin{tabularx}{\textwidth}{@{}>{\raggedright\arraybackslash}p{2.8cm}XXX@{}}
    \toprule
    \textbf{Level} & \textbf{The adversary can\ldots} & \textbf{Real-world example} & \textbf{Where we meet it} \\
    \midrule
    1. Eavesdropping only & see the ciphertexts that travel on the channel & someone sniffing the traffic of a public Wi-Fi network & perfect secrecy (\Cref{sec:perfect-secrecy}) \\
    2. Chosen plaintext & obtain the encryption of messages of their choice & making the victim encrypt an email written by the attacker & chosen-plaintext security \\
    3. Chosen ciphertext & obtain the decryption of ciphertexts of their choice & sending manipulated ciphertexts to a server and observing how it reacts & chosen-ciphertext security \\
    4. Active attacker & modify, inject and replay messages & a man-in-the-middle between a user and a website & integrity: MACs and digital signatures \\
    \bottomrule
  \end{tabularx}
  \caption{The capabilities of the adversary, from the weakest to the strongest.}
\end{table}

\textbf{The more power we give to the adversary, the stronger the guarantee} we obtain if the scheme still resists. Careful: it is not the scheme that becomes stronger, it is the \textit{guarantee} that becomes broader.

\begin{framedex}{The bridge test}
  If you test a bridge with a 40-ton truck and the bridge holds, you know it is safe for every lighter truck as well. The test did not make the bridge stronger: it gave you a guarantee that covers more situations.

  In the same way, if a scheme resists an adversary who can make it encrypt and decrypt whatever she wants, it also resists an adversary who only eavesdrops. Those extra capabilities are not useless for attacking, on the contrary they are usually very powerful weapons: it is a well-designed scheme that makes them useless.
\end{framedex}

\begin{framedex}{The same scheme in different threat models}
  The One-Time Pad (\Cref{sec:otp}) can be secure in one threat model and insecure in another:
  \begin{itemize}
    \item against an unbounded adversary that sees \textbf{one single} ciphertext $\to$ perfectly secure;
    \item if the key is \textbf{reused} and the adversary sees two ciphertexts $\to$ broken, because $c_1 \oplus c_2 = m_1 \oplus m_2$;
    \item against an \textbf{active} adversary $\to$ no integrity, because the adversary can flip bits of the message at will.
  \end{itemize}
  So the correct sentence is never ``the One-Time Pad is secure'', but \textbf{``the One-Time Pad is perfectly secret against an unbounded adversary that sees a single ciphertext''}. A security claim always comes with its threat model.
\end{framedex}

\subsection{Two families of security}

\begin{itemize}
  \item \textbf{Information-theoretic (unconditional) security} $\to$ the adversary is unbounded, $\varepsilon = 0$ (or extremely small), and \textbf{no assumptions} are needed. Very strong but, as Shannon's theorem will show, very expensive.
  \item \textbf{Computational (conditional) security} $\to$ the adversary runs in probabilistic polynomial time in the security parameter $n$ (the length of the key), and $\varepsilon = \negl(n)$. The proof relies on a \textbf{hardness assumption}, for example that factoring large numbers is hard.
\end{itemize}

In the computational case the proof is a \textbf{reduction}: ``if an adversary could break $\Pi$, then we could use it to solve a problem that we believe to be hard''. Since we believe the problem is hard, no such adversary can exist.

\subsection{From the template to rigorous definitions}\label{sec:rigorous}

The template of \Cref{obs:template} tells us \textit{how} a security definition is built. To make it rigorous we must write down, precisely, \textbf{the game} and \textbf{the class of adversaries}. Let's do it for the goal we study first: \textbf{confidentiality against an eavesdropper}, in both families.

\paragraph{Step 1: the game, written formally.} In cryptography a game is called an \textbf{experiment}. It is exactly the ``attack or retreat'' quiz: the experiment outputs $1$ when Eve guesses which of the two messages was encrypted.

\begin{frameddefn}[label={def:privk}]{The eavesdropping experiment $\PrivK_{\Adv,\Pi}$}
  Given a scheme $\Pi$ and an adversary $\Adv$:
  \begin{enumerate}
    \item $\Adv$ outputs two messages $m_0, m_1 \in \Mspace$ (in the computational setting they must have the same length, $|m_0| = |m_1|$);
    \item a key $k$ is generated ($k \sample \Kspace$, or $k \leftarrow \Gen(1^n)$ in the computational setting) and a uniform bit $b \sample \{0,1\}$ is chosen; $\Adv$ receives the \textbf{challenge ciphertext} $c = \Enc(k, m_b)$;
    \item $\Adv$ outputs a bit $b'$;
    \item the output of the experiment is $1$ if $b' = b$ (the adversary guessed correctly) and $0$ otherwise.
  \end{enumerate}
\end{frameddefn}

\paragraph{Step 2a: the information-theoretic version.} The class of adversaries is ``every algorithm'', with no limits, and the tolerance is zero.

\begin{frameddefn}[label={def:perfect-indist}]{Perfect indistinguishability}
  $\Pi$ is \textbf{perfectly indistinguishable} if for \textbf{every} adversary $\Adv$, even unbounded:
  \[
    \Pr\big[\PrivK_{\Adv,\Pi} = 1\big] = \frac{1}{2}
  \]
\end{frameddefn}

This definition is equivalent to perfect secrecy (Katz and Lindell, Lemma 2.6), as we discuss in \Cref{sec:going-further}.

\paragraph{Step 2b: the computational version.} Two new ingredients are needed:
\begin{itemize}
  \item \textbf{The security parameter} $n$ $\to$ the length of the key. The scheme becomes $\Pi = (\Gen, \Enc, \Dec)$, where $\Gen(1^n)$ generates a key of length $n$.
  \item \textbf{Negligible functions} $\to$ the precise meaning of ``a tiny probability''.
\end{itemize}

\begin{frameddefn}[label={def:negl}]{Negligible function}
  A function $f : \N \to \R_{\geq 0}$ is \textbf{negligible} if for every polynomial $p$ there exists $N$ such that for every $n > N$:
  \[
    f(n) < \frac{1}{p(n)}
  \]
  In words: $f$ goes to zero faster than the inverse of \textbf{any} polynomial. For example, $2^{-n}$ is negligible, while $1/n^{10}$ is not.
\end{frameddefn}

\begin{frameddefn}[label={def:eav}]{Computational security against an eavesdropper}
  $\Pi = (\Gen, \Enc, \Dec)$ has \textbf{indistinguishable encryptions in the presence of an eavesdropper} if \textbf{for every} probabilistic polynomial-time adversary $\Adv$ \textbf{there exists} a negligible function $\negl$ such that, for every $n$:
  \[
    \Pr\big[\PrivK_{\Adv,\Pi}(n) = 1\big] \leq \frac{1}{2} + \negl(n)
  \]
\end{frameddefn}

What the rigorous versions add to the template:
\begin{itemize}
  \item \textbf{The game is fully specified} $\to$ step by step, instead of a generic ``the adversary wins''.
  \item \textbf{The threat model is a precise class of algorithms} $\to$ ``every adversary'' or ``every probabilistic polynomial-time adversary''.
  \item \textbf{The tolerance is a function of $n$, not a fixed number} $\to$ and \textbf{the order of the quantifiers matters}: \textit{for every} adversary \textit{there exists} a negligible function, so the function may depend on the adversary.
  \item \textbf{No absolute value} $\to$ if an adversary guessed less than half of the times, flipping its answer would give an adversary that guesses more than half of the times; so the upper bound $\frac{1}{2} + \negl(n)$ is enough.
\end{itemize}

\section{The pillars of security}\label{sec:pillars}

When we say that a communication system is \textbf{secure}, we mean that it guarantees one or more of the following \textbf{security properties}. Each property is a distinct \textbf{security goal}: it has its own formal definition, it is threatened by a different kind of adversary, and it is achieved by different cryptographic tools.

\begin{frameddefn}{The three pillars of security}
  \begin{itemize}
    \item \textbf{Confidentiality} $\to$ only the intended receiver can learn the message: the ciphertext reveals \textbf{nothing} about the plaintext (at most its length).
    \item \textbf{Integrity} $\to$ any modification of the message during transmission is \textbf{detected} by the receiver.
    \item \textbf{Authenticity} $\to$ the receiver can verify that the message \textbf{really comes from the claimed sender} (data-origin authentication).
  \end{itemize}
\end{frameddefn}

\begin{table}[H]
  \centering
  \small
  \begin{tabularx}{\textwidth}{@{}lXX@{}}
    \toprule
    \textbf{Pillar} & \textbf{Who threatens it} & \textbf{Cryptographic tool} \\
    \midrule
    Confidentiality & \textit{passive} adversary (Eve): reads the channel & encryption (symmetric-key and public-key) \\
    Integrity & \textit{active} adversary (Mallory): modifies, injects, replays & message authentication codes (MACs), digital signatures \\
    Authenticity & \textit{active} adversary: impersonates the sender, forges messages & MACs (symmetric setting), digital signatures (asymmetric setting) \\
    \bottomrule
  \end{tabularx}
\end{table}

We wrote ``one or more'' on purpose: a scheme does not have to guarantee all three properties, and usually it does not. An encryption scheme gives only confidentiality, while a MAC gives integrity and authenticity but hides nothing.

\begin{framedobs}{}
  \begin{itemize}
    \item \textbf{Confidentiality does not imply integrity.} The One-Time Pad is perfectly secret, but anyone can modify its ciphertext in a controlled way. The pillars must be guaranteed \textbf{separately}, and a real-world scheme combines them (this is called \textit{authenticated encryption}).
    \item \textbf{In cryptography integrity and authenticity come together.} If a MAC or a digital signature verifies under Alice's key, it proves both that the message was not modified and that it was produced by Alice.
    \item Digital signatures also give \textbf{non-repudiation}: Alice cannot later deny having signed a message.
    \item The classic \textbf{CIA triad} of information security is \textit{Confidentiality, Integrity, Availability}. Availability (the service keeps working, for example against Denial of Service attacks) is a property of the whole system, and cryptography alone cannot guarantee it. That is why, in a cryptography course, \textbf{authenticity} takes its place.
    \item A related notion is \textbf{entity authentication}: proving \textit{who you are} in real time, not just who wrote a message. We will see it later with identification schemes.
  \end{itemize}
\end{framedobs}
